% Reset figure and table counters so appendix items are numbered A1, A2, ...
\renewcommand{\thefigure}{A\arabic{figure}}
\renewcommand{\thetable}{A\arabic{table}}
\setcounter{figure}{0}
\setcounter{table}{0}

% -----------------------------------------------------------------------
\section{Supplementary Figures}
\label{sec:app_figures}

\subsection{Preprocessing}

\begin{figure}[htbp]
    \centering
    \includegraphics[width=\textwidth]{figures/mad_distributions.pdf}
    \caption{MAD distributions for each omics view after ingestion. The dashed red
    line marks the MAD of the last retained feature; all features to the right were
    selected. Retained counts: 5,000 transcripts, 5,000 proteins, 6,000
    phosphosites.}
    \label{fig:app_mad}
\end{figure}

\clearpage
\subsection{MOFA model fitting}

\begin{figure}[htbp]
    \centering
    \includegraphics[width=0.85\textwidth]{figures/elbo_convergence.pdf}
    \caption{Final ELBO per random seed. The model with the highest ELBO (seed 2,
    $-3.72 \times 10^6$) was selected for all downstream analyses.}
    \label{fig:app_elbo}
\end{figure}

\begin{figure}[htbp]
    \centering
    \includegraphics[width=\textwidth]{figures/factor_violin_by_tissue.pdf}
    \caption{Distribution of MOFA+ factor scores split by tissue type (tumour vs.\
    normal) for all 13 retained factors. Factor~1 shows complete separation; no
    other factor reaches significance after Bonferroni correction.}
    \label{fig:app_violin}
\end{figure}

\clearpage
\subsection{Tumourigenesis trajectories}

\begin{figure}[htbp]
    \centering
    \includegraphics[width=0.85\textwidth]{figures/trajectory_arrows_umap.pdf}
    \caption{Per-patient tumourigenesis trajectories in the UMAP of MOFA+ factor
    scores (all 13 factors, all 204 samples). Each arrow connects the normal to the
    tumour sample of one patient.}
    \label{fig:app_traj_umap}
\end{figure}

\begin{figure}[htbp]
    \centering
    \includegraphics[width=0.85\textwidth]{figures/trajectory_arrows_pca_no_factor1.pdf}
    \caption{Trajectory arrows in PCA of factor scores with Factor~1 excluded.
    Projecting out the dominant tumour-normal axis reveals substructure in the
    remaining factors that is obscured in Figure~\ref{fig:traj_pca}.}
    \label{fig:app_traj_pca_nof1}
\end{figure}

\begin{figure}[htbp]
    \centering
    \includegraphics[width=0.85\textwidth]{figures/trajectory_arrows_umap_no_factor1.pdf}
    \caption{Trajectory arrows in UMAP of factor scores with Factor~1 excluded.
    Arrow directions are more heterogeneous than in the full-factor UMAP
    (Figure~\ref{fig:app_traj_umap}), confirming that Factor~1 drives the dominant
    shared trajectory.}
    \label{fig:app_traj_umap_nof1}
\end{figure}

\clearpage
\subsection{$\Delta Z$ heterogeneity}

\begin{figure}[htbp]
    \centering
    \includegraphics[width=\textwidth]{figures/delta_distributions.pdf}
    \caption{Per-factor $\Delta Z$ distributions across all 102 patients. Factor~1
    is strongly negative in all patients (universal tumourigenesis shift); the
    remaining factors show patient-level variation of varying magnitude.}
    \label{fig:delta_distributions}
\end{figure}

\begin{figure}[htbp]
    \centering
    \includegraphics[width=0.7\textwidth]{figures/magnitude_by_stage.pdf}
    \caption{L2 norm of $\Delta Z$ (transformation magnitude) stratified by tumour
    stage. No significant difference was observed (Kruskal-Wallis $H = 3.40$,
    $p = 0.183$).}
    \label{fig:mag_by_stage}
\end{figure}

\begin{figure}[htbp]
    \centering
    \includegraphics[width=\textwidth]{figures/repr_comparison_panel.pdf}
    \caption{2$\times$2 comparison of $\Delta Z$ representations: raw vs.\
    L2-normalised $\times$ PCA vs.\ UMAP, using all 13 factors. Points are
    coloured by transformation magnitude (viridis). The panel was used to select
    L2-normalised $\Delta Z$ as the primary clustering representation.}
    \label{fig:app_repr_panel}
\end{figure}

\clearpage
\subsection{Clustering}

\begin{figure}[htbp]
    \centering
    \includegraphics[width=0.6\textwidth]{figures/cluster_pca.pdf}
    \caption{PCA of L2-normalised $\Delta Z$ coloured by transformation subtype.
    The three subtypes overlap substantially in 2D, consistent with the UMAP in
    Figure~\ref{fig:cluster_umap} and the low silhouette scores.}
    \label{fig:app_cluster_pca}
\end{figure}

\begin{figure}[htbp]
    \centering
    \includegraphics[width=\textwidth]{figures/cluster_centroids.pdf}
    \caption{Factor-level centroid profiles for each transformation subtype in raw
    $\Delta Z$ space. All subtypes have a strongly negative Factor~1 component;
    differences in Factors 2--13 define the subtype-specific patterns.}
    \label{fig:app_centroids}
\end{figure}

\begin{figure}[htbp]
    \centering
    \includegraphics[width=\textwidth]{figures/dispersion_summary.pdf}
    \caption{Within-cluster dispersion (left) and pairwise centroid distances
    (right) in L2-normalised $\Delta Z$ space. Centroid distances (0.56--0.83)
    exceed mean within-cluster radii (0.48--0.54) for all pairs, indicating that
    the subtypes are more separated than they are internally dispersed.}
    \label{fig:app_dispersion}
\end{figure}

\clearpage
\subsection{Biological interpretation}

\begin{figure}[htbp]
    \centering
    \includegraphics[width=\textwidth]{figures/volcano_panel.pdf}
    \caption{Volcano plots of one-vs-rest $\Delta$Expression for each transformation
    subtype. Each point is a gene; $x$-axis shows median difference in
    $\Delta\text{Expr}$ between the focal subtype and all other patients;
    $y$-axis shows $-\log_{10}(\text{FDR})$. Top differentially expressed genes
    are labelled.}
    \label{fig:app_volcano}
\end{figure}

\clearpage
\subsection{Clinical associations}

\begin{figure}[htbp]
    \centering
    \includegraphics[width=0.7\textwidth]{figures/magnitude_by_cluster.pdf}
    \caption{Transformation magnitude (L2 norm of $\Delta Z$) per subtype.
    No significant difference was observed across subtypes
    (Kruskal-Wallis $H = 3.72$, $p = 0.156$).}
    \label{fig:app_mag_cluster}
\end{figure}

\begin{figure}[htbp]
    \centering
    \includegraphics[width=\textwidth]{figures/clinical_overview_heatmap.pdf}
    \caption{Descriptive clinical characteristics per transformation subtype.
    Three panels show demographics (sex, age, smoking), tumour stage distribution,
    and driver mutation prevalence. Values are percentages; median age is reported
    separately in Table~\ref{tab:app_clinical}.}
    \label{fig:app_clinical_overview}
\end{figure}

% -----------------------------------------------------------------------
\clearpage
\section{Supplementary Tables}
\label{sec:app_tables}

\subsection{Factor tissue separation statistics}

\begin{table}[htbp]
    \centering
    \caption{Paired Wilcoxon test for tumour-vs-normal separation per MOFA+ factor,
    applied to $\Delta Z$ values across all 102 patients. $p$-values are corrected
    by Benjamini-Hochberg (BH). Only Factor~1 is significant after correction.}
    \label{tab:app_tissue_stats}
    \begin{tabular}{lrrrr}
        \toprule
        Factor & $W$ & Raw $p$ & Rank-biserial $r$ & BH-adjusted $p$ \\
        \midrule
        Factor1  &     0 & $1.8 \times 10^{-18}$ & 1.000 & $2.4 \times 10^{-17}$ \\
        Factor2  &  2174 & 0.131                 & 0.172 & 0.394 \\
        Factor3  &  2284 & 0.253                 & 0.130 & 0.470 \\
        Factor4  &  2190 & 0.145                 & 0.166 & 0.394 \\
        Factor5  &  2453 & 0.562                 & 0.066 & 0.812 \\
        Factor6  &  2160 & 0.119                 & 0.178 & 0.394 \\
        Factor7  &  2480 & 0.625                 & 0.056 & 0.812 \\
        Factor8  &  2600 & 0.930                 & 0.010 & 0.930 \\
        Factor9  &  2197 & 0.152                 & 0.164 & 0.394 \\
        Factor10 &  2370 & 0.392                 & 0.098 & 0.637 \\
        Factor11 &  2576 & 0.866                 & 0.019 & 0.930 \\
        Factor12 &  2241 & 0.198                 & 0.147 & 0.429 \\
        Factor13 &  2521 & 0.725                 & 0.040 & 0.856 \\
        \bottomrule
    \end{tabular}
\end{table}

\subsection{Clinical characteristics per subtype}

\begin{table}[htbp]
    \centering
    \caption{Clinical characteristics per transformation subtype. Values are
    percentages except median age. Smoking: ever-smoker (current or reformed).
    Cluster 0 has 34 patients with available survival data (3 missing).}
    \label{tab:app_clinical}
    \begin{tabular}{lrrr}
        \toprule
        & Subtype 0 ($n = 37$) & Subtype 1 ($n = 25$) & Subtype 2 ($n = 40$) \\
        \midrule
        Median age (years)      & 66   & 58   & 65.5 \\
        Female (\%)             & 54.1 & 12.0 & 35.0 \\
        Ever-smoker (\%)        & 54.1 & 68.0 & 55.0 \\
        \midrule
        Stage I (\%)            & 56.8 & 52.0 & 52.5 \\
        Stage II (\%)           & 32.4 & 20.0 & 30.0 \\
        Stage III (\%)          & 10.8 & 28.0 & 17.5 \\
        \midrule
        \textit{KRAS} mut.\ (\%)   & 16.2 & 40.0 & 37.5 \\
        \textit{EGFR} mut.\ (\%)   & 56.8 & 16.0 & 37.5 \\
        \textit{TP53} mut.\ (\%)   & 43.2 & 64.0 & 70.0 \\
        ALK fusion (\%)            &  5.4 & 36.0 & 30.0 \\
        \bottomrule
    \end{tabular}
\end{table}

\subsection{Kaplan-Meier at-risk table}

\begin{table}[htbp]
    \centering
    \caption{Kaplan-Meier at-risk table. At-risk counts at each time point;
    cumulative censored and event counts are shown for each subtype. Subtype 0
    has 34 patients with survival data available.}
    \label{tab:app_km_atrisk}
    \setlength{\tabcolsep}{5pt}
    \begin{tabular}{r rrr rrr rrr}
        \toprule
        & \multicolumn{3}{c}{Subtype 0} & \multicolumn{3}{c}{Subtype 1} & \multicolumn{3}{c}{Subtype 2} \\
        \cmidrule(lr){2-4} \cmidrule(lr){5-7} \cmidrule(lr){8-10}
        Days & At risk & Censored & Events & At risk & Censored & Events & At risk & Censored & Events \\
        \midrule
           0 & 34 &  0 &  0 & 25 &  0 &  0 & 40 &  0 &  0 \\
         250 & 30 &  3 &  1 & 23 &  0 &  2 & 36 &  3 &  1 \\
         500 & 23 &  9 &  2 &  9 & 12 &  4 & 17 & 19 &  4 \\
         750 & 21 & 10 &  3 &  6 & 15 &  4 & 13 & 21 &  6 \\
        1000 & 17 & 11 &  6 &  3 & 17 &  5 & 11 & 22 &  7 \\
        1250 & 10 & 18 &  6 &  2 & 18 &  5 &  7 & 25 &  8 \\
        1500 &  8 & 20 &  6 &  1 & 19 &  5 &  6 & 26 &  8 \\
        1750 &  4 & 23 &  7 &  1 & 19 &  5 &  4 & 28 &  8 \\
        2000 &  0 & 27 &  7 &  0 & 20 &  5 &  0 & 32 &  8 \\
        \bottomrule
    \end{tabular}
\end{table}
